% =============================================================================
%  DỰ ÁN TỐT NGHIỆP — BÁO CÁO CHÍNH (BẢN OUTLINE SỰ KIỆN)
%  Đề tài: Phân tích Hiệu suất và Dự báo Sản lượng Hệ thống Điện Mặt Trời
%  Nhóm: The Outliers
%  Trường: Cao đẳng FPT Polytechnic — Chuyên ngành Phân tích Dữ liệu
% =============================================================================

\documentclass[12pt, a4paper]{report}

% ======================== PACKAGES ==========================================
\usepackage[utf8]{inputenc}
\usepackage[T1]{fontenc}
\usepackage[vietnamese]{babel}
\usepackage{geometry}
\geometry{top=2.5cm, bottom=2.5cm, left=3cm, right=2cm}
% ======================== CẤU HÌNH CODE PYTHON (MUTED STYLE) ==============================
\usepackage{listings}
\usepackage{xcolor} % Bắt buộc phải có để dùng \definecolor

% Định nghĩa lại bảng màu dịu, ít sặc sỡ và chuyên nghiệp hơn
\definecolor{mutedgreen}{rgb}{0.2,0.5,0.3}    % Màu xanh lá trầm cho comment
\definecolor{mutedblue}{rgb}{0.15,0.35,0.6}   % Màu xanh dương dịu cho keywords thay cho magenta
\definecolor{mutedgray}{rgb}{0.45,0.45,0.45}  % Màu xám cho số dòng
\definecolor{mutedpurple}{rgb}{0.45,0.2,0.55} % Màu tím trầm cho strings thay cho codepurple
\definecolor{softbackground}{rgb}{0.98,0.98,0.98} % Màu nền xám trắng siêu nhẹ tránh mỏi mắt
\definecolor{bordergray}{rgb}{0.85,0.85,0.85}  % Màu viền xám nhạt

\lstdefinestyle{pythonstyle}{
    backgroundcolor   = \color{softbackground},
    commentstyle      = \color{mutedgreen}\itshape, % Comment chữ nghiêng xanh trầm
    keywordstyle      = \color{mutedblue}\bfseries, % Keyword chữ đậm xanh dịu
    numberstyle       = \tiny\color{mutedgray},
    stringstyle       = \color{mutedpurple},
    basicstyle        = \ttfamily\small\color{black}, % Chữ code mặc định màu đen thuần
    breakatwhitespace = false,
    breaklines        = true,
    captionpos        = t, 
    keepspaces        = true,
    numbers           = left,
    numbersep         = 5pt,
    showspaces        = false,
    showstringspaces  = false,
    showtabs          = false,
    tabsize           = 4,
    frame             = single,
    rulecolor         = \color{bordergray}, % Viền mảnh màu xám nhạt tinh tế
    extendedchars     = true,
    % Bảng ánh xạ ký tự tiếng Việt (literate) đã làm sạch để fix lỗi UTF-8
    literate=
    {á}{{\'a}}1 {à}{{\`a}}1 {ã}{{\~a}}1 {ạ}{{\d{a}}}1 {ả}{{\h{a}}}1
    {ă}{{\u{a}}}1 {ắ}{{\'\u{a}}}1 {ằ}{{\`\u{a}}}1 {ẵ}{{\~\u{a}}}1 {ặ}{{\d{\u{a}}}}1 {ẳ}{{\h{\u{a}}}}1
    {â}{{\^a}}1 {ấ}{{\'\^a}}1 {ầ}{{\`\^a}}1 {ẫ}{{\~\^a}}1 {ậ}{{\d{\^a}}}1 {ẩ}{{\h{\^a}}}1
    {đ}{{\dj}}1 {Đ}{{\DJ}}1
    {é}{{\'e}}1 {è}{{\`e}}1 {ẽ}{{\~e}}1 {ẹ}{{\d{e}}}1 {ẻ}{{\h{e}}}1
    {ê}{{\^e}}1 {ế}{{\'\^e}}1 {ề}{{\`\^e}}1 {ễ}{{\~\^e}}1 {ệ}{{\d{\^e}}}1 {ể}{{\h{\^e}}}1
    {í}{{\'i}}1 {ì}{{\`i}}1 {ĩ}{{\~i}}1 {ị}{{\d{i}}}1 {ỉ}{{\h{i}}}1
    {ó}{{\'o}}1 {ò}{{\`o}}1 {õ}{{\~o}}1 {ọ}{{\d{o}}}1 {ỏ}{{\h{o}}}1
    {ô}{{\^o}}1 {ố}{{\'\^o}}1 {ồ}{{\`\^o}}1 {ỗ}{{\~\^o}}1 {ộ}{{\d{\^o}}}1 {ổ}{{\h{\^o}}}1
    {ơ}{{\v{o}}}1 {ớ}{{\'\v{o}}}1 {ờ}{{\`\v{o}}}1 {ỡ}{{\~\v{o}}}1 {ợ}{{\d{\v{o}}}}1 {ở}{{\h{\v{o}}}}1
    {ú}{{\'u}}1 {ù}{{\`u}}1 {ũ}{{\~u}}1 {ụ}{{\d{u}}}1 {ủ}{{\h{u}}}1
    {ư}{{\v{u}}}1 {ứ}{{\'\v{u}}}1 {ừ}{{\`\v{u}}}1 {ữ}{{\~\v{u}}}1 {ự}{{\d{\v{u}}}}1 {ử}{{\h{\v{u}}}}1
    {ý}{{\'y}}1 {ỳ}{{\`y}}1 {ỹ}{{\~y}}1 {ỵ}{{\d{y}}}1 {ỷ}{{\h{y}}}1
}
\lstset{style=pythonstyle}

% Font - Helvetica (Sans-Serif, Modern, supports Vietnamese)
\usepackage[scaled]{helvet}
\renewcommand{\familydefault}{\sfdefault}

% Header / Footer
\usepackage{fancyhdr}
\usepackage{lastpage}

% Hyperlinks & Bookmarks
\usepackage{hyperref}
\hypersetup{
    colorlinks  = true,
    linkcolor   = black,
    citecolor   = blue,
    filecolor   = magenta,
    urlcolor    = blue,
    pdftitle    = {Dự án Tốt nghiệp - Phân tích Hiệu suất và Dự báo Sản lượng Hệ thống Điện Mặt Trời},
    pdfauthor   = {The Outliers - FPT Polytechnic},
    bookmarksnumbered = true,
}

% Toán, ký hiệu
\usepackage{amsmath, amssymb}

% Danh sách
\usepackage{enumitem}
\usepackage{array}

% Bảng biểu
\usepackage{booktabs}
\usepackage{tabularx}
\usepackage{longtable}
\usepackage{multirow}
\usepackage{array}

% Hình ảnh
\usepackage{graphicx}
\usepackage{float}
\usepackage{caption}
\usepackage{subcaption}

% Màu sắc & Code
\usepackage{xcolor}
\usepackage{listings}

% Kẻ khung trang bìa
\usepackage{tikz}

% Danh mục hình, bảng
\usepackage[titles]{tocloft}

% Indent đoạn đầu
\usepackage{indentfirst}
\setlength{\parindent}{1cm}
\setlength{\parskip}{6pt}

% Line spacing
\usepackage{setspace}
\onehalfspacing
\setlength{\headheight}{33pt}

% ======================== CẤU HÌNH CODE PYTHON ==============================
\definecolor{codegreen}{rgb}{0,0.6,0}
\definecolor{codegray}{rgb}{0.5,0.5,0.5}
\definecolor{codepurple}{rgb}{0.58,0,0.82}
\definecolor{backcolour}{rgb}{0.96,0.96,0.94}
\definecolor{fptorange}{RGB}{245,130,32}

\lstdefinestyle{pythonstyle}{
    backgroundcolor   = \color{backcolour},
    commentstyle      = \color{codegreen},
    keywordstyle      = \color{magenta},
    numberstyle       = \tiny\color{codegray},
    stringstyle       = \color{codepurple},
    basicstyle        = \ttfamily\small,
    breakatwhitespace = false,
    breaklines        = true,
    captionpos        = r,
    keepspaces        = true,
    numbers           = left,
    numbersep         = 5pt,
    showspaces        = false,
    showstringspaces  = false,
    showtabs          = false,
    tabsize           = 4,
    frame             = single,
    rulecolor         = \color{codegray},
}
\lstset{style=pythonstyle}

% --- CẤU HÌNH KHOẢNG TRỐNG TIÊU ĐỀ CHƯƠNG ---
\usepackage{titlesec}
% Cấu hình cho \chapter và \chapter*
\titleformat{\chapter}[display]
  {\normalfont\huge\bfseries}{}{0pt}{\huge}
\titlespacing*{\chapter}{0pt}{-10pt}{20pt} 
% Tham số: {lề trái}{khoảng trống phía trên}{khoảng trống phía dưới tiêu đề}
% ======================== HEADER / FOOTER ====================================
% Trang nội dung chính (Logo ở header trái, THE OUTLIERS ở header phải)
\fancypagestyle{main}{
    \fancyhf{}
    \fancyhead[L]{\includegraphics[height=1cm]{images/logo_fpoly.png}}
    \fancyhead[R]{\small THE OUTLIERS}
    \fancyfoot[L]{\small Dự án Tốt nghiệp}
    \fancyfoot[R]{\small Page \thepage\ of \pageref{LastPage}}
    \renewcommand{\headrulewidth}{0.4pt}
    \renewcommand{\footrulewidth}{0.4pt}
}

% Trang đầu chương (plain)
\fancypagestyle{plain}{
    \fancyhf{}
    \fancyfoot[L]{\small Dự án Tốt nghiệp}
    \fancyfoot[R]{\small Page \thepage\ of \pageref{LastPage}}
    \renewcommand{\headrulewidth}{0pt}
    \renewcommand{\footrulewidth}{0.4pt}
}

% Trang mở đầu (roman numbering)
\fancypagestyle{front}{
    \fancyhf{}
    \fancyfoot[L]{\small Dự án Tốt nghiệp}
    \fancyfoot[R]{\small Page \thepage\ of \pageref{LastPage}}
    \renewcommand{\headrulewidth}{0pt}
    \renewcommand{\footrulewidth}{0.4pt}
}

% =============================================================================
%                           BẮT ĐẦU TÀI LIỆU
% =============================================================================
\begin{document}

% ────────────────────────────────────────────────────────────────────
%  TRANG BÌA (COVER PAGE)
% ────────────────────────────────────────────────────────────────────
\begin{titlepage}
\enlargethispage{1.2cm} % SỬA ĐỔI: Chỉ nới rộng tối đa 1.2cm để giữ chữ luôn nằm TRONG khung TikZ

\begin{tikzpicture}[remember picture, overlay]
    \draw[line width=2pt, fptorange]
        ([shift={(1cm,-1cm)}]current page.north west)
        rectangle
        ([shift={(-1cm,1cm)}]current page.south east);
    \draw[line width=0.5pt, fptorange]
        ([shift={(1.15cm,-1.15cm)}]current page.north west)
        rectangle
        ([shift={(-1.15cm,1.15cm)}]current page.south east);
\end{tikzpicture}

\begin{center}
    \textbf{\large TRƯỜNG CAO ĐẲNG FPT POLYTECHNIC}

    \vspace{0.5cm} % Thu nhỏ khoảng cách dưới tên trường
    % --- Logo trường Cao đẳng FPT Polytechnic ---
    \includegraphics[width=0.38\textwidth]{images/logo_fpoly.png}

    \vspace{0.4cm}
    {\Large \textbf{DỰ ÁN TỐT NGHIỆP}} \\[4pt]
    {\large \textbf{CHUYÊN NGÀNH: XỬ LÝ DỮ LIỆU}}

    \vspace{1.6cm}
    
    \begin{spacing}{1.0}
        {\huge \textbf{PHÂN TÍCH HIỆU SUẤT VÀ \\[4pt] DỰ BÁO SẢN LƯỢNG \\[4pt] HỆ THỐNG ĐIỆN MẶT TRỜI}}
    \end{spacing}

    \vspace{1.8cm} % Thu nhỏ khoảng cách trước khi vào phần thông tin nhóm
    \normalsize
    \begin{flushleft}
        \hspace{2cm} \textbf{Nhóm thực hiện:} The Outliers \\[5pt]
        \hspace{2cm} \textbf{Sinh viên thực hiện:} \\[3pt]
        % Tạo bảng phụ có chiều rộng cố định (12.5cm) để căn lề MSSV thẳng hàng bên phải
        \hspace{2cm}
        \begin{tabular*}{7.5cm}{@{} l @{\extracolsep{\fill}} r @{}}
            \hspace{0.5cm} 1. Nguyễn Vũ Nhựt   & PS44442 \\
            \hspace{0.5cm} 2. Nguyễn Văn Sỹ    & PS43671 \\
            \hspace{0.5cm} 3. Lê Công Toàn    & PS44145 \\
            \hspace{0.5cm} 4. Nguyễn Xuân Hùng & PS44527 \\
            \hspace{0.5cm} 5. Ngô Tấn Đạt     & PS44589 \\
        \end{tabular*} \\[5pt]
        \hspace{2cm} \textbf{Giảng viên hướng dẫn:} Văn Công Khanh
    \end{flushleft}
    \vfill
    {\normalsize \textbf{TP. HỒ CHÍ MINH, THÁNG 6 NĂM 2026}}
    \vspace{0.5cm} % Tạo một khoảng đệm nhỏ an toàn với đường viền dưới
\end{center}
\end{titlepage}

\pagenumbering{arabic}
\pagestyle{front}

% ────────────────────────────────────────────────────────────────────
%  LỜI CẢM ƠN
% ────────────────────────────────────────────────────────────────────
\chapter*{Lời cảm ơn}
\addcontentsline{toc}{chapter}{Lời cảm ơn}
\textit{[Nội dung lời cảm ơn của sinh viên đối với giảng viên hướng dẫn, gia đình và nhà trường sẽ được viết tại đây.]}

% ────────────────────────────────────────────────────────────────────
%  NHẬN XÉT CỦA GIÁO VIÊN
% ────────────────────────────────────────────────────────────────────
\chapter*{Nhận xét của Giảng viên Hướng dẫn}
\addcontentsline{toc}{chapter}{Nhận xét của Giảng viên Hướng dẫn}
\textit{[Nội dung nhận xét và đánh giá của Giảng viên Hướng dẫn đối với kết quả thực hiện dự án của nhóm.]}

\chapter*{Nhận xét của Giảng viên Phản biện}
\addcontentsline{toc}{chapter}{Nhận xét của Giảng viên Phản biện}
\textit{[Nội dung nhận xét và đánh giá của Giảng viên Phản biện tại buổi bảo vệ dự án tốt nghiệp.]}

% ────────────────────────────────────────────────────────────────────
%  TÓM TẮT DỰ ÁN
% ────────────────────────────────────────────────────────────────────
\chapter*{Tóm tắt Dự án}
\addcontentsline{toc}{chapter}{Tóm tắt Dự án}
Dự án tốt nghiệp ``Phân tích và Xử lý Bất thường trong Dữ liệu Sản lượng Năng lượng Mặt trời'' nhằm mục đích xây dựng một hệ thống phân tích dữ liệu toàn diện để phát hiện, xử lý và dự báo các bất thường trong sản lượng phát điện của các hệ thống quang điện (PV). 
Bằng cách thu thập và tích hợp dữ liệu từ 42 trạm phát điện tại Úc cùng dữ liệu khí tượng viễn thám từ Open-Meteo, nhóm đã thiết kế một kho dữ liệu đa chiều (Data Warehouse) trên nền tảng Supabase theo mô hình Galaxy Schema. Quy trình ETL được tự động hóa để làm sạch dữ liệu, xử lý các hiện tượng khuyết thiếu, nhiễu ban đêm, và sai lệch do cảm biến. 

Qua quá trình Phân tích Khám phá Dữ liệu (EDA), nhóm đã rút ra các insight quan trọng về sự ảnh hưởng của nhiệt độ, mức độ mây che phủ đến hiệu suất PV. Ngoài ra, dự án cũng áp dụng các mô hình học máy chuỗi thời gian như ARIMA và Prophet để thiết lập đường cơ sở (Baseline) dự báo sản lượng, kết hợp hiển thị trực quan thông qua Dashboard Tableau. Hệ thống không chỉ cung cấp cái nhìn toàn cảnh về tình trạng hoạt động của các trạm điện mặt trời, mà còn hỗ trợ ban quản lý đưa ra quyết định kịp thời nhằm tối ưu hóa hiệu suất và giảm thiểu rủi ro vận hành.

% ────────────────────────────────────────────────────────────────────
%  MỤC LỤC & DANH MỤC
% ────────────────────────────────────────────────────────────────────
\tableofcontents
\listoffigures
\listoftables

\chapter*{Danh mục Từ viết tắt}
\addcontentsline{toc}{chapter}{Danh mục Từ viết tắt}
\vspace*{-15pt}

\noindent\hspace*{-6pt}%
    \begin{tabularx}{\textwidth}{@{} l X @{}}
    \toprule
    \textbf{Ký hiệu/Từ viết tắt} & \textbf{Ý nghĩa đầy đủ} \\
    \midrule
    API     & Application Programming Interface (Giao diện lập trình ứng dụng) \\
    ARIMA   & AutoRegressive Integrated Moving Average \\
    AWS     & Amazon Web Services \\
    BI      & Business Intelligence (Phân tích thông minh doanh nghiệp) \\
    CF      & Capacity Factor (Hệ số công suất) \\
    CSDL    & Cơ sở dữ liệu \\
    CT      & Current Transformer (Cảm biến dòng điện) \\
    DWH     & Data Warehouse (Kho dữ liệu) \\
    EDA     & Exploratory Data Analysis (Phân tích khám phá dữ liệu) \\
    ETL     & Extract, Transform, Load (Trích xuất, Biến đổi, Nạp) \\
    FK      & Foreign Key (Khóa ngoại) \\
    GCP     & Google Cloud Platform \\
    IoT     & Internet of Things (Internet vạn vật) \\
    IQR     & Interquartile Range (Khoảng tứ phân vị) \\
    kWp     & Kilowatt-peak (Công suất đỉnh cực đại) \\
    LSTM    & Long Short-Term Memory \\
    PK      & Primary Key (Khóa chính) \\
    PR      & Performance Ratio (Hệ số hiệu suất) \\
    PV      & Photovoltaic (Quang điện - Năng lượng mặt trời) \\
    UI/UX   & User Interface / User Experience (Giao diện / Trải nghiệm người dùng) \\
    WMO     & World Meteorological Organization (Tổ chức Khí tượng Thế giới) \\
    XGBoost & Extreme Gradient Boosting \\
    \bottomrule
    \end{tabularx}

\newpage
\pagestyle{main}

% ====================================================================
%  CHƯƠNG 1: TỔNG QUAN ĐỀ TÀI VÀ BÀI TOÁN NGHIÊN CỨU
% ====================================================================
\chapter{Tổng quan Đề tài và Bài toán Nghiên cứu}
\label{chap:tongquan}

\section{Bối cảnh thực tế và lý do chọn đề tài}
\label{sec:lydo}

Năng lượng mặt trời đang ngày càng đóng vai trò quan trọng trong cơ cấu năng lượng tái tạo toàn cầu. Theo Cơ quan Năng lượng Quốc tế (IEA), công suất lắp đặt điện mặt trời trên toàn thế giới đã vượt mốc 1.600~GW vào cuối năm 2023, với tốc độ tăng trưởng trung bình 25\%/năm trong thập kỷ qua. Tại Việt Nam, chính sách khuyến khích phát triển điện mặt trời áp mái và điện mặt trời tập trung đã đưa tổng công suất lắp đặt lên hơn 16.500~MW. Tuy nhiên, hiệu suất thực tế của các hệ thống quang điện (PV) thường sai lệch đáng kể so với giá trị thiết kế do ảnh hưởng của nhiều yếu tố như: điều kiện thời tiết biến đổi, suy giảm chất lượng tấm pin, bụi bẩn, bóng che và lỗi thiết bị.

Bài toán phân tích hiệu suất và phát hiện sớm các bất thường trong vận hành là then chốt giúp tối ưu hoá hiệu quả đầu tư và nâng cao độ tin cậy của lưới điện. Lý do nhóm chọn đề tài này bao gồm:
\begin{itemize}[noitemsep]
    \item \textbf{Tính thực tiễn:} Dự báo sản lượng và phát hiện bất thường có ứng dụng trực tiếp cho các doanh nghiệp vận hành nhà máy điện mặt trời.
    \item \textbf{Tính liên ngành:} Kết hợp phân tích dữ liệu, thống kê, học máy và kiến trúc kho dữ liệu trong một đề tài thống nhất.
    \item \textbf{Nguồn dữ liệu mở:} Dữ liệu từ Open-Meteo API miễn phí, chất lượng cao và phủ rộng toàn cầu.
    \item \textbf{So sánh đa vùng:} Phân tích đồng thời dữ liệu từ nhiều khu vực tại Úc để đánh giá tổng quát về hiệu suất PV.
\end{itemize}

\section{Bài toán kinh doanh và các khó khăn cần giải quyết}
\label{sec:problem}

Trong vận hành thực tế các trạm điện mặt trời, ban quản lý thường đối mặt với các khó khăn lớn:
\begin{itemize}[noitemsep]
    \item \textbf{Tác động thời tiết bất định:} Sự thay đổi của mây che phủ, bức xạ và nhiệt độ môi trường khiến sản lượng điện biến động liên tục, gây khó khăn cho việc điều phối lưới điện.
    \item \textbf{Suy hao và hỏng hóc thiết bị:} Các tấm pin bị bám bụi, inverter bị lỗi phần sụn hoặc quá nhiệt làm suy giảm đáng kể hiệu suất phát điện mà không được phát hiện kịp thời.
    \item \textbf{Sai lệch số liệu giám sát:} Cảm biến đo bức xạ hoặc cảm biến CT đo dòng điện bị nhiễu xung điện đột ngột hoặc lệch pha tần suất dữ liệu (Granularity Mismatch) giữa dữ liệu sản lượng (15 phút) và dữ liệu thời tiết (1 giờ) khiến phân tích sai lệch.
\end{itemize}
Hệ quả của việc không giải quyết các vấn đề trên là doanh nghiệp ra quyết định chậm trễ, tổn thất tài chính lớn do mất sản lượng, và dữ liệu thu thập bị bẩn dẫn đến các mô hình dự báo không chính xác.

\section{Mục tiêu dự án và Kết quả dự kiến}
\label{sec:muctieu}

Dự án tốt nghiệp đặt ra các mục tiêu cốt lõi sau:
\begin{enumerate}[noitemsep]
    \item Thiết kế và triển khai kho dữ liệu (\textit{Data Warehouse}) tích hợp đồng bộ dữ liệu sản lượng và thời tiết theo mô hình Star Schema trên nền tảng Supabase.
    \item Xây dựng pipeline dữ liệu tự động làm sạch (xử lý khuyết thiếu, lọc nhiễu ban đêm, phát hiện bất thường cảm biến).
    \item Thực hiện phân tích khám phá chuyên sâu (EDA) để rút ra các insight vận hành quan trọng như hiện tượng suy giảm hiệu suất do nhiệt độ và độ nhiễu ban đêm.
    \item Huấn luyện các mô hình dự báo sản lượng Baseline (ARIMA, Prophet) và so sánh hiệu năng.
\end{enumerate}
\textbf{Kết quả dự kiến (Sneak Peek):} Một hệ thống Dashboard Tableau trực quan hóa mượt mà cho phép các nhà quản lý theo dõi toàn diện hiệu suất vận hành của 42 trạm tại Úc, tự động cảnh báo lỗi cảm biến và hiển thị dự báo sản lượng cho các chu kỳ tiếp theo với độ tin cậy cao.

\section{Đối tượng và phạm vi nghiên cứu}
\label{sec:phamvi}

\begin{itemize}[noitemsep]
    \item \textbf{Đối tượng nghiên cứu:} Hiệu suất phát điện của hệ thống quang điện (PV) và các yếu tố khí tượng viễn thám (bức xạ sóng ngắn, nhiệt độ, độ phủ mây, lượng mưa).
    \item \textbf{Phạm vi không gian:} 42 trạm điện mặt trời tại Úc và 63 tỉnh thành Việt Nam.
    \item \textbf{Phạm vi thời gian:} Dữ liệu lịch sử từ 01/01/2020 đến 30/04/2022.
\end{itemize}

\section{Bố cục báo cáo}
\label{sec:bocuc}

Báo cáo dự án tốt nghiệp được cấu trúc thành 7 chương chính như sau:
\begin{itemize}[noitemsep]
    \item \textbf{Chương 1:} Tổng quan Đề tài và Bài toán Nghiên cứu.
    \item \textbf{Chương 2:} Cơ sở Lý thuyết (các mô hình dự báo, phát hiện outlier và DWH).
    \item \textbf{Chương 3:} Thu thập Dữ liệu và Thiết kế Kho dữ liệu.
    \item \textbf{Chương 4:} Quy trình ETL, Làm sạch và Biến đổi Dữ liệu.
    \item \textbf{Chương 5:} Phân tích Khám phá Dữ liệu (EDA) và Trực quan hóa.
    \item \textbf{Chương 6:} Mô hình hóa và Dự báo Sản lượng Baseline.
    \item \textbf{Chương 7:} Đánh giá Dự án, Đề xuất và Theo dõi Vận hành.
\end{itemize}


% ====================================================================
%  CHƯƠNG 2: CƠ SỞ LÝ THUYẾT
% ====================================================================
\chapter{Cơ sở Lý thuyết}
\label{chap:lythuyet}

\section{Tổng quan về Năng lượng mặt trời và Bức xạ quang điện}
\label{sec:solar_theory}

Hệ thống quang điện (PV) chuyển đổi năng lượng bức xạ mặt trời trực tiếp thành điện năng xoay chiều thông qua hiệu ứng quang điện trong chất bán dẫn. Hiệu suất chuyển đổi phụ thuộc chặt chẽ vào các chỉ số:
\begin{itemize}[noitemsep]
    \item \textbf{Performance Ratio (PR):} Tỷ lệ giữa sản lượng thực tế tạo ra và sản lượng lý thuyết trong điều kiện tối ưu.
    \item \textbf{Capacity Factor (CF):} Tỷ lệ giữa sản lượng thực tế và sản lượng tối đa trạm có thể phát nếu chạy liên tục ở công suất thiết kế cực đại.
    \item \textbf{Specific Yield:} Sản lượng điện tạo ra trên mỗi đơn vị công suất đỉnh lắp đặt (kWh/kWp).
\end{itemize}

\section{Kiến trúc Kho dữ liệu đa chiều (Data Warehouse)}
\label{sec:dwh_theory}

Kiến trúc kho dữ liệu đa chiều theo phương pháp Kimball tập trung thiết kế các bảng theo hai nhóm thực thể:
\begin{itemize}[noitemsep]
    \item \textbf{Bảng Fact (Sự kiện):} Chứa các số liệu đo lường liên tục thay đổi (Measures) và các khóa ngoại trỏ tới các bảng chiều.
    \item \textbf{Bảng Dimension (Chiều):} Chứa các thông tin mô tả ngữ cảnh nền tảng (Thời gian, vị trí địa lý, thông tin thiết bị) được chuẩn hóa để làm bộ lọc chéo dữ liệu.
\end{itemize}
Mô hình Galaxy Schema (Fact Constellation) là sự mở rộng của Star Schema, cho phép nhiều bảng Fact dùng chung các bảng chiều (Conformed Dimensions) để giải quyết các dữ liệu có tần suất khác nhau (Granularity Mismatch).

\section{Quy trình Trích xuất, Biến đổi và Nạp dữ liệu (ETL/ELT Pipeline)}
\label{sec:etl_theory}

Quy trình ETL bao gồm ba bước cơ bản:
\begin{enumerate}[noitemsep]
    \item \textbf{Extract (Trích xuất):} Đọc dữ liệu thô từ các nguồn phân tán (API, database vật lý, file CSV).
    \item \textbf{Transform (Biến đổi):} Làm sạch, chuẩn hóa kiểu dữ liệu, nội suy các bản ghi khuyết thiếu và xử lý giá trị ngoại lai.
    \item \textbf{Load (Nạp):} Đẩy dữ liệu sạch vào Kho dữ liệu tập trung (Data Warehouse) để phục vụ phân tích.
\end{enumerate}

\section{Các phương pháp Phát hiện bất thường dữ liệu}
\label{sec:anomaly_theory}

\begin{itemize}[noitemsep]
    \item \textbf{Phương pháp thống kê Z-Score:} Xác định outlier dựa trên độ lệch chuẩn so với giá trị trung bình:
    \begin{equation}
        Z = \frac{x - \mu}{\sigma}
    \end{equation}
    Bản ghi bị gắn cờ outlier nếu $|Z| > 3$.
    \item \textbf{Phương pháp khoảng tứ phân vị IQR:} Outlier là điểm nằm ngoài biên dưới hoặc biên trên:
    \begin{equation}
        [\text{Lower Bound}, \text{Upper Bound}] = [Q_1 - 1.5 \cdot \text{IQR}, Q_3 + 1.5 \cdot \text{IQR}]
    \end{equation}
    trong đó $\text{IQR} = Q_3 - Q_1$.
    \item \textbf{Phân tích phần dư (Residual Analysis):} Đánh giá độ chênh lệch tuyệt đối $|e_t| = |y_t - \hat{y}_t|$ giữa giá trị thực tế và giá trị dự báo từ mô hình học máy. Các điểm có phần dư vượt quá ngưỡng sai số quy định sẽ bị xem là bất thường.
\end{itemize}

\section{Cơ sở thuật toán Dự báo chuỗi thời gian}
\label{sec:forecast_theory}

\begin{itemize}[noitemsep]
    \item \textbf{Mô hình ARIMA:} Mô hình AutoRegressive Integrated Moving Average kết hợp các thành phần tự hồi quy (p), sai phân (d) và trung bình trượt (q):
    \begin{equation}
        \phi(B)(1-B)^d X_t = \theta(B) \varepsilon_t
    \end{equation}
    \item \textbf{Mô hình Prophet:} Mô hình phân tách chuỗi thời gian do Meta phát triển, xử lý tốt tính mùa vụ mạnh và các giá trị khuyết thiếu:
    \begin{equation}
        y(t) = g(t) + s(t) + h(t) + \varepsilon_t
    \end{equation}
    với $g(t)$ đại diện cho xu hướng, $s(t)$ cho tính mùa vụ, $h(t)$ cho ngày lễ và $\varepsilon_t$ cho sai số ngẫu nhiên.
\end{itemize}


% ====================================================================
%  CHƯƠNG 3: THU THẬP DỮ LIỆU VÀ THIẾT KẾ KHO DỮ LIỆU
% ====================================================================
\chapter{Thu thập Dữ liệu và Thiết kế Kho dữ liệu}
\label{chap:thuthap_kho}

\section{Phương pháp thu thập dữ liệu thô (Extract)}
\label{sec:datannguon}

Dự án thu thập dữ liệu từ hai nguồn chính:
\begin{enumerate}[noitemsep]
    \item \textbf{Dữ liệu giám sát trạm phát điện thực tế:} Ghi nhận sản lượng điện của 42 trạm tại Úc dưới dạng tệp tin CSV cục bộ và các thông tin cấu hình trạm.
    \item \textbf{Dữ liệu khí tượng viễn thám:} Thu thập từ Open-Meteo API. Quy trình gọi API sử dụng tọa độ địa lý (vĩ độ, kinh độ) của các trạm để tải về chuỗi thời gian khí tượng theo giờ bao gồm bức xạ sóng ngắn, nhiệt độ, lượng mưa và mây che phủ.
\end{enumerate}

\subsection{Trích xuất dữ liệu khí tượng qua Open-Meteo API}
Dữ liệu khí tượng viễn thám được thu thập tự động thông qua Open-Meteo Archive API. Quy trình sử dụng thư viện \texttt{requests} và \texttt{pandas} để đọc tọa độ (vĩ độ, kinh độ) của 42 trạm phát từ tệp danh mục \texttt{Solar\_Site\_Details.csv}. Các biến số quan trọng như bức xạ sóng ngắn (\texttt{shortwave\_radiation}), nhiệt độ (\texttt{temperature\_2m}), mã thời tiết (\texttt{weather\_code}), và trạng thái ngày/đêm (\texttt{is\_day}) được trích xuất theo chu kỳ giờ.

Để đảm bảo tính ổn định khi gọi API hàng loạt, mã nguồn được tích hợp cơ chế tự động xử lý lỗi vượt quá giới hạn truy cập (Rate Limit - HTTP 429). Khi máy chủ trả về mã 429, hệ thống sẽ tạm dừng (\texttt{time.sleep}) 60 giây trước khi thử lại, giúp luồng trích xuất không bị gián đoạn.

\begin{lstlisting}[language=Python, caption={Mã nguồn Python trích xuất dữ liệu khí tượng từ Open-Meteo API}, label={lst:api_weather}]
import requests
import pandas as pd
import time

# 1. Khoi tao cau hinh tham so API
START_DATE = "2023-01-01"
END_DATE = "2023-12-31"
FINAL_OUTPUT_FILE = "open_meteo_weather_raw_2023.csv"

HOURLY_COLUMNS = [
    "shortwave_radiation", "direct_radiation", 
    "diffuse_radiation", "temperature_2m", 
    "weather_code", "is_day", "cloud_cover",
    "wind_speed_10m", "precipitation", "sunshine_duration"
]

# Doc toa do tu danh muc tram
site_df = pd.read_csv('/content/Solar_Site_Details(1).csv')
url = "https://archive-api.open-meteo.com/v1/archive"
all_data_frames = []

# 2. Vong lap lay du lieu va xu ly Rate Limit
for index, row in site_df.iterrows():
    params = {
        "latitude": row['lat'],
        "longitude": row['Lon'],
        "start_date": START_DATE,
        "end_date": END_DATE,
        "hourly": ",".join(HOURLY_COLUMNS),
        "timezone": "Asia/Ho_Chi_Minh"
    }
    
    success = False
    while not success:
        try:
            response = requests.get(
            url, 
            params=params, 
            timeout=120)
            
            if response.status_code == 200:
                data = response.json()
                if "hourly" in data:
                    df_site = pd.DataFrame(
                    {"timestamp": data["hourly"]["time"]}
                    )
                    for col in HOURLY_COLUMNS:
                        df_site[col] = data["hourly"].get(col)
                    
                    df_site["SiteKey"] = row.get('SiteKey', index)
                    all_data_frames.append(df_site)
                    success = True
            elif response.status_code == 429:
                print("Rate limited. Waiting 60s...")
                time.sleep(60)
            else:
                success = True # Bo qua neu loi dinh dang khac
        except Exception as e:
            time.sleep(5)
            success = True

# Xuat tep du lieu tong hop
if all_data_frames:
    final_df = pd.concat(all_data_frames, ignore_index=True)
    final_df['timestamp'] = pd.to_datetime(final_df['timestamp'])
    final_df.to_csv(FINAL_OUTPUT_FILE, index=False)
\end{lstlisting}

\subsection{Thiết lập kết nối cơ sở dữ liệu Supabase}
Để lưu trữ và xử lý tập dữ liệu sau khi trích xuất, dự án cấu hình hệ quản trị cơ sở dữ liệu PostgreSQL trên nền tảng Supabase (Khu vực Southeast Asia). Điểm đặc thù trong kiến trúc kết nối và nạp (Load) dữ liệu của nhóm:
\begin{itemize}[noitemsep]
    \item \textbf{Sử dụng Connection Pooler:} Do Supabase free tier không hỗ trợ kết nối trực tiếp qua IPv4, dự án cấu hình kết nối thông qua Supabase Pooler \\(host \texttt{aws-1-ap-southeast-1.pooler.supabase.com}, cổng 5432). Giao thức này hỗ trợ IPv4 ổn định cho môi trường làm việc nhóm, đồng thời quản lý và tái sử dụng các phiên kết nối để tránh tình trạng quá tải (Connection Overload).
    \item \textbf{Tối ưu hóa Thư viện kết nối (\texttt{pg8000}):} Thay vì sử dụng thư viện \texttt{psycopg2} truyền thống, dự án chuyển sang sử dụng \texttt{pg8000}. Đây là thư viện Pure Python, giúp khắc phục triệt để lỗi thiếu hụt các thư viện hệ thống (\texttt{libz.so.1}, \texttt{libstdc++.so.6}) trên môi trường NixOS, đảm bảo mã nguồn có thể chạy mượt mà trên mọi hệ điều hành.
    \item \textbf{Bảo mật hệ thống:} Toàn bộ thông tin nhạy cảm bao gồm Mật khẩu cơ sở dữ liệu và Project ID được mã hóa và truyền qua tệp \texttt{.env} thay vì khai báo trực tiếp trong mã nguồn. Mã nguồn nạp dữ liệu vào kho được xử lý cách ly, cam kết không rò rỉ (commit) mật khẩu lên nền tảng GitHub.
\end{itemize}
Bộ dữ liệu thô được lưu trữ thành các tệp tin CSV: \texttt{campus\_meta.csv}, \texttt{Solar\_Site\_Details.csv},\\ \texttt{calender.csv}, \texttt{Solar\_Energy\_Generation.csv}, và \texttt{open\_meteo\_weather\_raw\_2023.csv}.

% ====================================================================
% MỤC MỚI BỔ SUNG: TỪ ĐIỂN DỮ LIỆU & META DATA
% ====================================================================
\section{Data Dictionary \& Metadata}
\label{sec:data_dictionary}

Sau quá trình trích xuất và thu thập, bộ dữ liệu thô được cô đọng lại thành 5 tệp tin cốt lõi (Core Datasets). Việc xây dựng từ điển dữ liệu (Data Dictionary) ở giai đoạn này giúp định nghĩa rõ ràng cấu trúc, kiểu dữ liệu và ý nghĩa của từng trường thông tin gốc trước khi thực hiện các bước biến đổi (Transform) và đưa vào Kho dữ liệu.

\subsection{Dữ liệu Thông tin Cơ sở hạ tầng (\texttt{campus\_meta.csv})}
Tệp lưu trữ thông tin định danh cấp cao nhất về các khu vực/cơ sở dự án. Quy mô: 5 dòng.
\begin{table}[H]
\centering
\caption{Từ điển dữ liệu tệp \texttt{campus\_meta.csv}}
\renewcommand{\arraystretch}{1.4}
\begin{tabularx}{\textwidth}{p{3.5cm} c p{2cm} X}
\toprule
\textbf{Tên trường} & \textbf{Kiểu DL} & \textbf{Đơn vị} & \textbf{Ý nghĩa \& Mô tả kỹ thuật} \\
\midrule
id       & Integer & -  & Mã định danh duy nhất (Primary Key) cho từng cơ sở. \\
name     & VARCHAR  & -  & Tên gọi chính thức của cơ sở (VD: Bundoora, Bendigo...). \\
capacity & Integer & kW & Sức chứa, quy mô hoạt động hoặc mức công suất thiết kế. \\
\bottomrule
\end{tabularx}
\end{table}

\newpage
\subsection{Dữ liệu Thông tin Kỹ thuật Trạm (\texttt{Solar\_Site\_Details.csv})}
Tệp lưu trữ chi tiết cấu hình kỹ thuật, thiết bị lắp đặt và tọa độ địa lý của các trạm phát điện. Quy mô: 42 dòng.
\begin{table}[H]
\centering
\caption{Từ điển dữ liệu tệp \texttt{Solar\_Site\_Details.csv}}
\renewcommand{\arraystretch}{1.4}
\begin{tabularx}{\textwidth}{p{3.5cm} c p{2cm} X}
\toprule
\textbf{Tên trường} & \textbf{Kiểu DL} & \textbf{Đơn vị} & \textbf{Ý nghĩa \& Mô tả kỹ thuật} \\
\midrule
CampusKey        & Integer & -   & Mã định danh cơ sở nơi đặt trạm điện. \\
SiteKey          & Integer & -   & Mã định danh duy nhất của từng trạm phát điện. \\
kWp              & Float   & kWp & Công suất đỉnh tối đa của hệ thống. \\
Number of panels & Float   & Tấm & Tổng số lượng tấm pin mặt trời được lắp đặt. \\
Panel            & String  & -   & Thương hiệu và công suất định mức của tấm pin. \\
Inverter         & String  & -   & Thông tin bộ biến đổi dòng điện (Inverter). \\
Optimizers       & String  & -   & Thông tin bộ tối ưu hóa công suất đi kèm. \\
Metric           & String  & -   & Đơn vị đo lường sản lượng điện năng. \\
lat              & Float   & Độ ($^\circ$) & Vĩ độ địa lý (Latitude). \\
Lon              & Float   & Độ ($^\circ$) & Kinh độ địa lý (Longitude). \\
\bottomrule
\end{tabularx}
\end{table}

\subsection{Dữ liệu Sản lượng Điện (\texttt{Solar\_Energy\_Generation.csv})}
Bảng dữ liệu chuỗi thời gian (Fact Table) chứa sản lượng phát điện thực tế được ghi nhận liên tục theo chu kỳ 15 phút. Quy mô: 2.731.947 dòng.
\begin{table}[H]
\centering
\caption{Từ điển dữ liệu tệp \texttt{Solar\_Energy\_Generation.csv}}
\renewcommand{\arraystretch}{1.4}
\begin{tabularx}{\textwidth}{p{3.5cm} c p{2cm} X}
\toprule
\textbf{Tên trường} & \textbf{Kiểu DL} & \textbf{Đơn vị} & \textbf{Ý nghĩa \& Mô tả kỹ thuật} \\
\midrule
CampusKey       & Integer  & -   & Mã định danh cơ sở liên kết. \\
SiteKey         & Integer  & -   & Mã định danh trạm phát điện mặt trời cụ thể. \\
Timestamp       & DateTime & -   & Mốc thời gian đo đạc dữ liệu sản lượng (Chu kỳ 15 phút). \\
SolarGeneration & Float    & kWh & Sản lượng điện năng thực tế tạo ra trong khoảng 15 phút đó. \\
\bottomrule
\end{tabularx}
\end{table}

\subsection{Dữ liệu Thông tin Lịch trình (\texttt{calender.csv})}
Bảng dữ liệu đánh dấu các sự kiện thời gian theo cấp độ ngày, giúp phân tích các yếu tố hành vi ảnh hưởng đến nhu cầu tiêu thụ và vận hành. Quy mô: 2.313 dòng.
\begin{table}[H]
\centering
\caption{Từ điển dữ liệu tệp \texttt{calender.csv}}
\renewcommand{\arraystretch}{1.4}
\begin{tabularx}{\textwidth}{p{3.5cm} c p{2cm} X}
\toprule
\textbf{Tên trường} & \textbf{Kiểu DL} & \textbf{Đơn vị} & \textbf{Ý nghĩa \& Mô tả kỹ thuật} \\
\midrule
date         & Date    & - & Mốc thời gian theo ngày (định dạng DD/MM/YYYY). \\
is\_holiday  & Boolean & - & Cờ đánh dấu ngày nghỉ lễ (1: Lễ, 0: Bình thường). \\
is\_semester & Boolean & - & Cờ đánh dấu ngày nằm trong giai đoạn học kỳ chính thức. \\
is\_exam     & Boolean & - & Cờ đánh dấu ngày diễn ra kỳ thi tại cơ sở. \\
\bottomrule
\end{tabularx}
\end{table}

\subsection{Dữ liệu Khí tượng Viễn thám (\texttt{open\_meteo\_weather\_raw.csv})}
Chuỗi thời gian chứa các chỉ số khí tượng được trích xuất qua API tương ứng với tọa độ từng trạm phát. Quy mô: 367.920 dòng.
\begin{table}[H]
\centering
\caption{Từ điển dữ liệu tệp \texttt{open\_meteo\_weather\_raw.csv}}
\renewcommand{\arraystretch}{1.4}
\begin{tabularx}{\textwidth}{p{3.8cm} c p{2cm} X}
\toprule
\textbf{Tên trường} & \textbf{Kiểu DL} & \textbf{Đơn vị} & \textbf{Ý nghĩa \& Mô tả kỹ thuật} \\
\midrule
shortwave\_radiation & Float   & W/m$^2$ & Tổng bức xạ sóng ngắn bề mặt. Feature có tương quan cao nhất với sản lượng điện. \\
direct\_radiation    & Float   & W/m$^2$ & Bức xạ mặt trời chiếu trực tiếp tới bề mặt. \\
diffuse\_radiation   & Float   & W/m$^2$ & Bức xạ mặt trời bị khuếch tán (do mây, bụi, sương mù). \\
temperature\_2m      & Float   & $^\circ$C & Nhiệt độ không khí ở độ cao 2 mét. Phục vụ đánh giá suy giảm nhiệt. \\
weather\_code        & Integer & -       & Mã thời tiết theo chuẩn WMO (0: Quang đãng, 51-67: Mưa...). Phân loại Categorical tốt cho ML. \\
\bottomrule
\end{tabularx}
\end{table}
% ====================================================================
\section{Thiết kế kiến trúc Kho dữ liệu đa chiều (Data Warehouse)}
\label{sec:star_schema}

Quá trình thiết kế kho dữ liệu (Data Warehouse - DWH) được thực hiện theo phương pháp từ trên xuống (Top-down) kết hợp mô hình đa chiều của Kimball. Nhằm giải quyết bài toán dữ liệu lớn và độ lệch pha thời gian (Granularity Mismatch) giữa dữ liệu sản lượng (15 phút) và dữ liệu khí tượng (1 giờ), hệ thống được thiết kế theo cấu trúc \textbf{Lược đồ Thiên hà (Galaxy Schema / Fact Constellation)}. Quá trình này bao gồm 3 cấp độ mô hình hóa: Conceptual, Logical, và Physical.

\subsection{Mô hình Dữ liệu Khái niệm (Conceptual Data Model)}
Ở mức khái niệm, hệ thống xác định các thực thể kinh doanh cốt lõi và mối quan hệ giữa chúng mà không đi sâu vào chi tiết kỹ thuật. Mô hình khái niệm của hệ thống được phân rã thành các thực thể (như minh họa tại Hình \ref{fig:conceptual_model}):
\begin{itemize}[noitemsep]
    \item \textbf{Thực thể trung tâm (Đo lường năng lượng):} Là nơi hội tụ các chỉ số đo lường hiệu suất.
    \item \textbf{Các thực thể ngữ cảnh (Nguyên nhân/Điều kiện):} Bao gồm \textit{Trạm năng lượng (Solar site)}, \textit{Vị trí địa lý}, \textit{Ngày (Date)} và \textit{Thời tiết}.
\end{itemize}

\begin{figure}[H]
    \centering
    \includegraphics[width=0.85\textwidth]{images/Mo_hinh_khai_niem_by_Toan.png}
    \caption{Sơ đồ Mô hình Dữ liệu Khái niệm (Conceptual Data Model)}
    \label{fig:conceptual_model}
\end{figure}

Mối quan hệ cốt lõi: Vị trí địa lý, Thời tiết và Ngày có quan hệ 1-M (one-many) với Đo lường năng lượng, giúp ban quản lý hiểu rõ dữ liệu sản lượng và thời tiết dưới góc nhìn không gian và thời gian.

\subsection{Mô hình Dữ liệu Logic (Logical Data Model)}
Mô hình logic cụ thể hóa các thực thể thành các bảng và xác định các thuộc tính (attributes), khóa chính (Primary Key), khóa ngoại (Foreign Key). Dựa trên thiết kế kiến trúc (xem Hình \ref{fig:logical_model}), hệ thống gồm 5 bảng chiều và 2 bảng sự kiện:

\begin{enumerate}[label=\arabic*., noitemsep]
    \item \textbf{Các bảng chiều dùng chung (Conformed Dimensions):}
    \begin{itemize}[noitemsep]
        \item \texttt{dim\_solar\_site}: Khóa chính \texttt{site\_id}. Lưu thông tin thiết bị: \texttt{campus\_name}, \texttt{capacity\_kw}, \texttt{Number\_of\_panels}, \texttt{Panel}, \texttt{Inverter}, \texttt{Optimizers}, \texttt{Metric}.
        \item \texttt{dim\_geography}: Khóa chính \texttt{geo\_id}. Định vị không gian: \texttt{latitude}, \texttt{longitude}, \texttt{location\_name}.
        \item \texttt{dim\_date}: Khóa chính \texttt{date\_id}. Trục ngày: \texttt{full\_date}, \texttt{day}, \texttt{month}, \texttt{year}.
        \item \texttt{dim\_time}: Khóa chính \texttt{time\_id}. Trục giờ chi tiết (chu kỳ 15p): \texttt{time\_string}, \texttt{hour}, \texttt{minute}.
        \item \texttt{dim\_weather\_type}: Khóa chính \texttt{weather\_type\_id}. Phân loại trạng thái: \texttt{weather\_code}, \texttt{is\_day}, \texttt{weather\_condition}, \texttt{description}.
    \end{itemize}
    \item \textbf{Các bảng sự kiện (Fact Tables):}
    \begin{itemize}[noitemsep]
        \item \texttt{fact\_solar\_energy\_gen}: Khóa chính \texttt{gen\_id}. Kết nối với Site, Geo, Date, Time. Đo lường sản lượng \texttt{energy\_generated\_kwh}.
        \item \texttt{fact\_weather}: Khóa chính \texttt{weather\_id}. Kết nối với Geo, Date, Time, Weather\_type. Chứa các chỉ số viễn thám như \texttt{shortwave\_radiation}, \texttt{temperature\_c}, \texttt{cloud\_cover\_total}, \texttt{wind\_speed}, \texttt{precipitation\_mm}.
    \end{itemize}
\end{enumerate}

\begin{figure}[H]
    \centering
    \includegraphics[width=\textwidth]{images/Mo_hinh_logic_by_Hung.png}
    \caption{Sơ đồ Mô hình Dữ liệu Logic (Logical Data Model)}
    \label{fig:logical_model}
\end{figure}

\subsection{Mô hình Dữ liệu Vật lý (Physical Data Model)}
Mô hình vật lý định nghĩa cách thức dữ liệu được lưu trữ thực tế trên hệ quản trị cơ sở dữ liệu (sơ đồ ở Hình \ref{fig:physical_model}). Ở bước này, các kiểu dữ liệu (Data Types) cụ thể (INT, FLOAT, VARCHAR, DATE) và ràng buộc khóa ngoại (Foreign Keys) được thiết lập để tối ưu hóa không gian lưu trữ và tốc độ truy vấn. Chi tiết cấu trúc vật lý của các bảng chiều được mô tả từ Bảng \ref{tab:dim_solar_site} đến Bảng \ref{tab:dim_weather_type}. Các bảng sự kiện sản lượng và thời tiết được mô tả tương ứng trong Bảng \ref{tab:fact_gen_spec} và Bảng \ref{tab:fact_weather_spec}.

\begin{figure}[H]
    \centering
    \includegraphics[width=\textwidth]{images/Mo_hinh_Phisical_By_Toan.png}
    \caption{Sơ đồ Mô hình Dữ liệu Vật lý (Physical Data Model)}
    \label{fig:physical_model}
\end{figure}

% =====================================================================
% BẢNG 1: DIM SOLAR SITE
% =====================================================================
\begin{table}[H]
\centering
\caption{Cấu trúc vật lý bảng chiều trạm năng lượng (\texttt{dim\_solar\_site})}
\label{tab:dim_solar_site}
\renewcommand{\arraystretch}{1.4}
\begin{tabularx}{\textwidth}{p{4.5cm} c X}
\toprule
\textbf{Tên trường (Field)} & \textbf{Kiểu DL} & \textbf{Diễn giải vật lý} \\
\midrule
site\_id           & INT (PK) & Khóa chính của trạm. \\
campus\_name       & VARCHAR  & Tên khu vực đặt trạm dưới dạng chuỗi ký tự. \\
capacity\_kw       & FLOAT    & Công suất thiết kế (số thực). \\
Number\_of\_panels & INT      & Tổng số lượng tấm pin (số nguyên). \\
Panel              & VARCHAR  & Thương hiệu/loại tấm pin (chuỗi ký tự). \\
Inverter           & VARCHAR  & Thông tin bộ biến đổi điện (chuỗi ký tự). \\
Optimizers         & VARCHAR  & Bộ tối ưu hóa (chuỗi ký tự). \\
Metric             & VARCHAR  & Ghi chú/tiêu chuẩn đo lường (chuỗi ký tự). \\
\bottomrule
\end{tabularx}
\end{table}

% =====================================================================
% BẢNG 2: DIM GEOGRAPHY
% =====================================================================
\begin{table}[H]
\centering
\caption{Cấu trúc vật lý bảng chiều địa lý (\texttt{dim\_geography})}
\label{tab:dim_geography}
\renewcommand{\arraystretch}{1.4}
\begin{tabularx}{\textwidth}{p{4.5cm} c X}
\toprule
\textbf{Tên trường (Field)} & \textbf{Kiểu DL} & \textbf{Diễn giải vật lý} \\
\midrule
geo\_id        & INT (PK) & Khóa chính định danh vị trí địa lý. \\
latitude      & FLOAT    & Vĩ độ (số thực). \\
longitude     & FLOAT    & Kinh độ (số thực). \\
location\_name & VARCHAR  & Tên địa danh hành chính. \\
capacity      & INT      & Tổng công suất của khu vực đặt trạm (kW). \\
\bottomrule
\end{tabularx}
\end{table}

% =====================================================================
% BẢNG 3: DIM DATE
% =====================================================================
\begin{table}[H]
\centering
\caption{Cấu trúc vật lý bảng chiều ngày (\texttt{dim\_date})}
\label{tab:dim_date}
\renewcommand{\arraystretch}{1.4}
\begin{tabularx}{\textwidth}{p{4.5cm} c X}
\toprule
\textbf{Tên trường (Field)} & \textbf{Kiểu DL} & \textbf{Diễn giải vật lý} \\
\midrule
date\_id       & INT (PK) & Khóa chính đại diện ngày. \\
full\_date     & DATE     & Trường lưu định dạng thời gian nguyên bản (YYYY-MM-DD). \\
day            & INT      & Ngày trong tháng (số nguyên). \\
month          & INT      & Tháng trong năm (số nguyên). \\
year           & INT      & Năm phân tích (số nguyên). \\
is\_holiday    & INT      & Đánh dấu ngày lễ (1: Ngày lễ, 0: Ngày thường). \\
is\_semester   & INT      & Đánh dấu học kỳ (1: Trong kỳ học, 0: Kỳ nghỉ). \\
is\_exam       & INT      & Đánh dấu mùa thi (1: Trong mùa thi, 0: Bình thường). \\
\bottomrule
\end{tabularx}
\end{table}

% =====================================================================
% BẢNG 4: DIM TIME
% =====================================================================
\begin{table}[H]
\centering
\caption{Cấu trúc vật lý bảng chiều thời gian (\texttt{dim\_time})}
\label{tab:dim_time}
\renewcommand{\arraystretch}{1.4}
\begin{tabularx}{\textwidth}{p{4.5cm} c X}
\toprule
\textbf{Tên trường (Field)} & \textbf{Kiểu DL} & \textbf{Diễn giải vật lý} \\
\midrule
time\_id     & INT (PK) & Khóa chính mốc thời gian. \\
time\_string & VARCHAR  & Chuỗi ký tự biểu diễn giờ (HH:MM). \\
hour         & INT      & Giờ trong ngày (số nguyên). \\
minute       & INT      & Phút (số nguyên). \\
\bottomrule
\end{tabularx}
\end{table}

% =====================================================================
% BẢNG 5: DIM WEATHER TYPE
% =====================================================================
\begin{table}[H]
\centering
\caption{Cấu trúc vật lý bảng chiều danh mục thời tiết (\texttt{dim\_weather\_type})}
\label{tab:dim_weather_type}
\renewcommand{\arraystretch}{1.4}
\begin{tabularx}{\textwidth}{p{4.5cm} c X}
\toprule
\textbf{Tên trường (Field)} & \textbf{Kiểu DL} & \textbf{Diễn giải vật lý} \\
\midrule
weather\_type\_id    & INT (PK) & Khóa chính danh mục thời tiết. \\
weather\_code       & INT      & Mã trạng thái WMO (số nguyên). \\
is\_day             & INT      & Trạng thái ánh sáng (1: Ngày, 0: Đêm) lưu dạng số nguyên. \\
weather\_condition  & VARCHAR  & Tên điều kiện thời tiết tổng quát. \\
description         & VARCHAR  & Chuỗi văn bản mô tả chi tiết diễn biến thời tiết. \\
\bottomrule
\end{tabularx}
\end{table}

% =====================================================================
% BẢNG 6: FACT SOLAR ENERGY GEN
% =====================================================================
\begin{table}[H]
\centering
\caption{Cấu trúc vật lý bảng sự kiện sản lượng (\texttt{fact\_solar\_energy\_gen})}
\label{tab:fact_gen_spec}
\renewcommand{\arraystretch}{1.4}
\begin{tabularx}{\textwidth}{p{4.5cm} c X}
\toprule
\textbf{Tên trường (Field)} & \textbf{Kiểu DL} & \textbf{Diễn giải mục đích vật lý} \\
\midrule
gen\_id & INT & Khóa chính định danh từng dòng sự kiện (Auto-increment). \\
site\_id & INT & Khóa ngoại trỏ đến \texttt{dim\_solar\_site.site\_id}. \\
geo\_id & INT & Khóa ngoại trỏ đến \texttt{dim\_geography.geo\_id}. \\
date\_id & INT & Khóa ngoại trỏ đến \texttt{dim\_date.date\_id}. \\
time\_id & INT & Khóa ngoại trỏ đến \texttt{dim\_time.time\_id}. \\
energy\_generated\_kwh & FLOAT & Sản lượng điện năng sinh ra (kiểu số thực để tính toán độ chính xác cao). \\
\bottomrule
\end{tabularx}
\end{table}

% =====================================================================
% BẢNG 7: FACT WEATHER
% =====================================================================
\begin{table}[H]
\centering
\caption{Cấu trúc vật lý bảng sự kiện thời tiết (\texttt{fact\_weather})}
\label{tab:fact_weather_spec}
\renewcommand{\arraystretch}{1.4}
\begin{tabularx}{\textwidth}{p{5cm} c X}
\toprule
\textbf{Tên trường (Field)} & \textbf{Kiểu DL} & \textbf{Diễn giải mục đích vật lý} \\
\midrule
weather\_id & INT & Khóa chính định danh dòng lịch sử thời tiết. \\
geo\_id & INT & Khóa ngoại trỏ đến \texttt{dim\_geography.geo\_id}. \\
date\_id & INT & Khóa ngoại trỏ đến \texttt{dim\_date.date\_id}. \\
time\_id & INT & Khóa ngoại trỏ đến \texttt{dim\_time.time\_id}. \\
weather\_type\_id & INT & Khóa ngoại trỏ đến \texttt{dim\_weather\_type.weather\_type\_id}. \\
is\_day & INT & Chỉ số phân tách chu kỳ ánh sáng (số nguyên: 1/0). \\
shortwave\_radiation & INT & Bức xạ sóng ngắn (số nguyên). \\
temperature\_c & FLOAT & Nhiệt độ không khí (số thực). \\
cloud\_cover\_total & FLOAT & Tổng độ che phủ mây (số thực, tỷ lệ \%). \\
cloud\_cover\_low & FLOAT & Tỷ lệ mây tầng thấp (số thực). \\
cloud\_cover\_mid & FLOAT & Tỷ lệ mây tầng trung (số thực). \\
cloud\_cover\_high & FLOAT & Tỷ lệ mây tầng cao (số thực). \\
Diffuse\_Solar\_Radiation & INT & Ánh sáng tán xạ (số nguyên). \\
Direct\_Normal\_Irradiance & INT & Bức xạ chiếu thẳng (số nguyên). \\
wind\_speed & FLOAT & Tốc độ gió (số thực). \\
precipitation\_mm & FLOAT & Tổng lượng mưa tích tụ (số thực). \\
Sunshine\_Duration & FLOAT & Thời lượng nắng (số thực). \\
\bottomrule
\end{tabularx}
\end{table}

\section{Kiểm tra tính toàn vẹn và bảo mật dữ liệu}
\label{sec:data_integrity}
Để đảm bảo tính toàn vẹn dữ liệu (Data Integrity) trong CSDL vật lý, các ràng buộc tham chiếu (Referential Constraints) được cấu hình nghiêm ngặt theo quan hệ Một - Nhiều (1-M). Mọi bản ghi trong các bảng Fact bắt buộc phải có khóa tham chiếu tồn tại trong các bảng Dimension tương ứng, ngăn chặn tuyệt đối tình trạng dữ liệu mồ côi (orphan records) khi thực hiện luồng ETL.


% ====================================================================
%  CHƯƠNG 4: QUY TRÌNH ETL VÀ CHUẨN HÓA DỮ LIỆU
% ====================================================================
\chapter{Quy trình ETL và Chuẩn hóa Dữ liệu}
\label{chap:etl_cleaning}

\section{Tổng quan Kiến trúc Pipeline Xử lý Dữ liệu}
\label{sec:pipeline_overview}
\textit{[Mô tả luồng dữ liệu tự động (Data Flow) từ nguồn dữ liệu thô (API, CSV), qua các bước trích xuất, làm sạch bằng Python, và tải (Load) vào kho dữ liệu Supabase theo đúng chuẩn Star Schema.]}

\section{Đánh giá Chất lượng và Khám phá Dữ liệu (EDA ban đầu)}
\label{sec:data_quality}
\textit{[Kiểm tra dữ liệu khuyết thiếu, trùng lặp và các lỗi định dạng ban đầu của dữ liệu thời tiết và sản lượng.]}

\section{Làm sạch và Chuẩn hóa Dữ liệu (Data Cleaning)}
\label{sec:cleaning}
\subsection{Xử lý Dữ liệu Khuyết thiếu và Lỗi Logic}
\textit{[Sử dụng các phương pháp nội suy (Interpolation) để điền khuyết dữ liệu thời tiết. Đồng nhất các đơn vị đo lường và định dạng mốc thời gian (Timestamps).]}

\subsection{Xử lý Dữ liệu Ngoại lai (Outliers) và Lọc nhiễu}
\textit{[Áp dụng phương pháp khoảng tứ phân vị (IQR) để phát hiện và xử lý các bản ghi đột biến do lỗi cảm biến. Xây dựng thuật toán lọc nhiễu ban đêm (công suất phát ảo khi không có bức xạ).]}

\section{Biến đổi và Trích xuất Đặc trưng (Feature Engineering)}
\label{sec:transformation}
\textit{[Tạo lập các đặc trưng thời gian (Giờ, Ngày, Tháng, Mùa vụ). Gom cụm dữ liệu (Aggregation) theo chu kỳ 1 giờ để giải quyết bài toán lệch pha tần suất giữa dữ liệu thời tiết và sản lượng.]}


% ====================================================================
%  CHƯƠNG 5: TRỰC QUAN DỮ LIỆU VÀ XÂY DỰNG DASHBOARD
% ====================================================================
\chapter{Trực quan Dữ liệu và Xây dựng Dashboard}
\label{chap:dashboard}

\section{Chiến lược Kết nối Dữ liệu và Thiết kế UI/UX}
\label{sec:dashboard_strategy}
\textit{[Kết nối các bảng Fact và Dimension từ Supabase vào công cụ BI (Tableau) thông qua các Relationship (1-M). Lựa chọn bộ màu sắc và bố cục Dashboard theo quy luật Gestalt để làm nổi bật các chỉ số quan trọng.]}

\section{Chi tiết Hệ thống Dashboard Phân tích}
\label{sec:dashboard_details}
\subsection{Dashboard 1: Tổng quan Vận hành (Executive Overview)}
\textit{[Cung cấp cái nhìn toàn cảnh về sản lượng điện tổng, công suất hoạt động trung bình, và xu hướng theo thời gian của 42 trạm phát.]}

\subsection{Dashboard 2: Phân tích Hiệu suất và Khí hậu}
\textit{[Phân tích tương quan giữa sản lượng điện với nhiệt độ và bức xạ. Phát hiện các mốc suy giảm hiệu suất.]}

\subsection{Dashboard 3: Giám sát Bất thường và Lỗi Cảm biến}
\textit{[Trực quan hóa các điểm dữ liệu bất thường (Outliers), tần suất xuất hiện lỗi tại các trạm để kỹ sư có kế hoạch bảo trì.]}


% ====================================================================
%  CHƯƠNG 6: KHAI PHÁ DỮ LIỆU VÀ BUSINESS INSIGHTS
% ====================================================================
\chapter{Khai phá Dữ liệu và Business Insights}
\label{chap:insights}

Theo tinh thần phân tích chuyên sâu (Đừng chỉ mô tả $\rightarrow$ phải giải thích), chương này tập trung vào việc bóc tách các điểm nghẽn vận hành và đưa ra các phát hiện có giá trị kinh doanh (Business Insights) dựa trên logic dữ liệu, thay vì chỉ dựa vào các biểu đồ bề nổi. Các mô hình học máy cơ sở (ARIMA/Prophet) được tích hợp như một công cụ chẩn đoán bổ trợ.

\section{Insight 1: Hiện tượng suy giảm hiệu suất do nhiệt độ (Thermal Degradation)}
\textit{[Dữ liệu chứng minh rằng khi nhiệt độ môi trường vượt quá ngưỡng giới hạn (ví dụ: $25^\circ\text{C}$), hiệu suất chuyển đổi quang điện của tấm pin sụt giảm đáng kể. Điều này giải thích tại sao vào buổi trưa bức xạ cao nhất nhưng công suất lại không đạt đỉnh.]}

\section{Insight 2: Điểm mù vận hành từ ``Độ nhiễu ban đêm''}
\textit{[Phát hiện hàng loạt trạm ghi nhận sản lượng phát điện vào ban đêm. Nguyên nhân được phân tích là do nhiễu xung điện từ biến tần (Inverter) hoặc dòng rò. Nếu không lọc bỏ, dữ liệu này sẽ làm sai lệch nghiêm trọng các báo cáo tổng sản lượng.]}

\section{Insight 3: Ứng dụng Mô hình dự báo (Baseline) để tối ưu vận hành}
\textit{[Ứng dụng mô hình Time-Series (như Prophet) để thiết lập đường cơ sở (Baseline) dự báo sản lượng. Sự chênh lệch lớn giữa sản lượng thực tế và Baseline đóng vai trò là ``cờ báo'' (Alert) để ban quản lý cử đội bảo trì đến kiểm tra các tấm pin bị bụi bẩn bám dính hoặc hỏng hóc.]}


% ====================================================================
%  CHƯƠNG 7: KẾT LUẬN VÀ HƯỚNG PHÁT TRIỂN
% ====================================================================
\chapter{Kết luận và Hướng phát triển}
\label{chap:conclusion}

\section{Kết luận Tổng quan}
\label{sec:conclusion}
\textit{[Tóm tắt lại việc hoàn thành các mục tiêu dự án: Xây dựng thành công Data Warehouse, làm sạch dữ liệu nhiễu, và phát triển hệ thống Dashboard trực quan giúp doanh nghiệp chuyển đổi từ quản lý thụ động sang giám sát dựa trên dữ liệu.]}

\section{Khuyến nghị Chiến lược (Actionable Insights)}
\label{sec:recommendations}
\textit{[Đề xuất các hành động cụ thể: Lắp đặt hệ thống làm mát tự động để giảm suy hao nhiệt, thiết lập quy trình kiểm tra Inverter định kỳ để khắc phục độ nhiễu ban đêm, và ứng dụng kết quả dự báo để lập kế hoạch vệ sinh tấm pin.]}

\section{Hạn chế của Dự án}
\label{sec:limitations}
\textit{[Những giới hạn về phần cứng, thời gian xử lý dữ liệu lớn, hoặc sự thiếu hụt của các mô hình học sâu (Deep Learning) phức tạp hơn.]}

\section{Hướng phát triển trong tương lai}
\label{sec:future_work}
\subsection{Nâng cấp Hệ thống Phân tích Dự đoán (Predictive Analytics)}
\textit{[Tích hợp các mô hình Machine Learning mạnh mẽ hơn (XGBoost, LSTM) để dự báo biến động công suất theo thời gian thực.]}

\subsection{Tự động hóa và Đưa lên Đám mây (Cloud \& Automation)}
\textit{[Triển khai toàn bộ Pipeline ETL lên các dịch vụ Cloud như AWS/GCP để xử lý luồng dữ liệu Streaming từ các cảm biến IoT.]}

\newpage
% ====================================================================
%  PHỤ LỤC A: KẾ HẠCH THỰC HIỆN
% ====================================================================
\chapter*{Phụ lục A: Kế hoạch Thực hiện}
\addcontentsline{toc}{chapter}{Phụ lục A: Kế hoạch Thực hiện}
\label{app:kehoach}
\textit{[Bảng phân công nhiệm vụ chi tiết và kế hoạch 14 tuần thực hiện dự án của các thành viên.]}


% ====================================================================
%  PHỤ LỤC B: CẤU TRÚC MÃ NGUỒN
% ====================================================================
\chapter*{Phụ lục B: Cấu trúc Mã nguồn}
\addcontentsline{toc}{chapter}{Phụ lục B: Cấu trúc Mã nguồn}
\label{app:code}
\textit{[Sơ đồ cấu trúc thư mục của package Python du\_an\_tot\_nghiep và các script tiện ích.]}


% ====================================================================
%  TÀI LIỆU THAM KHẢO
% ====================================================================
\chapter*{Tài liệu Tham khảo}
\addcontentsline{toc}{chapter}{Tài liệu Tham khảo}
\label{chap:thamkhao}

\begin{enumerate}[label={[\arabic*]}, noitemsep]
    \item Box, G. E. P., Jenkins, G. M., Reinsel, G. C., \& Ljung, G. M. (2015). \textit{Time Series Analysis: Forecasting and Control}. 5th ed. Wiley.
    \item Taylor, S. J., \& Letham, B. (2018). Forecasting at Scale. \textit{The American Statistician}, 72(1), 37--45.
    \item Breiman, L. (2001). Random Forests. \textit{Machine Learning}, 45(1), 5--32.
    \item Chen, T., \& Guestrin, C. (2016). XGBoost: A Scalable Tree Boosting System. \textit{Proceedings of KDD 2016}, 785--794.
    \item Hochreiter, S., \& Schmidhuber, J. (1997). Long Short-Term Memory. \textit{Neural Computation}, 9(8), 1735--1780.
    \item Vaswani, A., et al. (2017). Attention Is All You Need. \textit{Proceedings of NeurIPS 2017}, 5998--6008.
    \item Kimball, R., \& Ross, M. (2013). \textit{The Data Warehouse Toolkit}. 3rd ed. Wiley.
    \item NASA POWER Project. (2024). \url{https://power.larc.nasa.gov/}.
    \item Supabase Documentation. (2024). \url{https://supabase.com/docs}.
    \item IEA -- International Energy Agency. (2024). \url{https://www.iea.org/reports/renewables-2024}.
    \item Chandola, V., Banerjee, A., \& Kumar, V. (2009). Anomaly Detection: A Survey. \textit{ACM Computing Surveys}, 41(3), Article 15.
    \item Tukey, J. W. (1977). \textit{Exploratory Data Analysis}. Addison-Wesley.
\end{enumerate}

\end{document}